\section{2026: The Second Rupture}

\subsection{Announcement and the February exchange}

At the beginning of 2026 the Society published its annual statistics as
of 1 November 2025: 1{,}482 members, of whom two bishops, 733 priests,
264 seminarians and postulants, 145 brothers, 88 oblates and 250 sisters,
of fifty nationalities, with an average age of forty-seven. The
publication does not name the two bishops; they were Bernard Fellay and
Alfonso de Galarreta, the two who would confer the consecrations of
1 July. These are the Society's own figures, published on its own
information service.\endnote{``SSPX Statistics 2025'' (fsspx.news, 4 January 2026),
read 2026-07-25. Self-reported; no independent verification exists or is
attempted here. The same page states that the Society's bishops ``neither
receive nor claim any episcopal jurisdiction over priests or the
faithful.''} A body with two bishops and 733 priests dependent on them for ordination
and confirmation faced, in 2026, exactly the problem the founder had
faced in 1988.

On 2 February 2026 the Society's General House announced that episcopal
consecrations would take place at Écône on 1 July 2026. On 12 February
the Superior General, Father Davide Pagliarani, met Cardinal Víctor
Manuel Fernández, Prefect of the Dicastery for the Doctrine of the Faith.
According to the Society's communiqué of 19 February, the cardinal
proposed ``a path of specifically theological dialogue, according to a
very precise methodology,'' to establish ``the minimum requirements for
full communion with the Catholic Church,'' on condition that the
announced consecrations be suspended; the Superior General submitted the
proposal to his council and answered in writing on 18 February, declining
to make the consecrations conditional on the dialogue.\endnote{``Communiqué
from the General House: the Society's response to Rome'' (fsspx.news,
19 February 2026), read 2026-07-25, with the accompanying dossier
``Episcopal Consecrations: Updated Dossier'' (fsspx.news, from 24 March
2026). Party sources reporting a private meeting and publishing the
Society's own correspondence; the quotations above are the Society's
English rendering of what it says the cardinal proposed. Bounded
negative: the same communiqué refers to ``the statement published by the
Holy See on 12 February,'' but no statement concerning the Society
appears among the items of the Holy See Press Office bulletins for
February 2026 (all bulletin listings for that month checked 2026-07-25).
Whatever was issued on 12 February was not released through the ordinary
bulletin channel, and this account does not reproduce it.}

\subsection{The Roman warnings of May and June}

The Holy See's first published act of 2026 in this matter is a
declaration of Cardinal Fernández, issued through the press bulletin on
13 May 2026. It is four sentences long and does nothing but restate
existing law. Concerning the Society, it says, what has already been
communicated is reaffirmed: the episcopal ordinations announced by the
Society do not have the corresponding pontifical mandate; this gesture
``will constitute an act of a schismatic nature,'' quoting \work{Ecclesia
Dei} n.~3; and ``formal adherence to the schism constitutes a grave
offence against God and carries the excommunication established by the
law of the Church,'' quoting \work{Ecclesia Dei} 5~c) and referring to
the 1996 explanatory note of the Pontifical Council for Legislative
Texts. The Holy Father, it adds, continues to ask the Holy Spirit to
enlighten those responsible so that they may turn back from the most
grave decision they have taken.\endnote{\work{Dichiarazione di Sua
Eminenza il Card.\ Víctor Manuel Fernández, Prefetto del Dicastero per la
Dottrina della Fede}, Holy See Press Office bulletin B0403 of 13 May 2026
(\url{https://press.vatican.va/content/salastampa/it/bollettino/pubblico/2026/05/13/0403/00788.html}),
read 2026-07-25. Italian original; the English above is a paraphrase with
the quoted phrases rendering the declaration's own quotations of
\work{Ecclesia Dei}.}

On 26 May the Society published the names of the four priests to be
consecrated: Pascal Schreiber, rector of the seminary at Zaitzkofen;
Michael Goldade, rector of the seminary at Dillwyn; Michel Poinsinet de
Sivry, superior of the Benelux district; and Marc Hanappier, professor at
Dillwyn. On 24 June it published an open letter to Leo XIV, signed by the
Superior General with his assistants and counsellors, accompanied by a
lengthy \textit{Profession of Catholic Faith}.\endnote{The names and
offices are given in the declaration read at the consecrations (below)
and in the Society's timeline of 14 July 2026; the open letter and
profession are listed and published in the Society's 2026 dossier. Party
sources throughout.}

On 29 June 2026, the solemnity of Saints Peter and Paul, Pope Leo XIV
wrote to the Superior General. The letter is short, and its official
English text is unambiguous. The Church, he writes, recognizes the
devotion to liturgical life, the commitment to priestly formation, the
apostolic zeal and the desire for fidelity to Tradition that characterize
many people and communities connected to the Society, and this has
motivated his predecessors' attentive and generous attitude. Then:
``filled with Christian affection, I plead with you and ask you with all
my heart: please turn back! I urge you to consider carefully the
spiritual good of the faithful, because the schismatic act you are about
to undertake would deprive them of the licit and, in some cases, even
valid reception of the Sacraments, which they love and seek for their
sanctification.'' The Church remains open to a path of dialogue; but ``to
tear the seamless garment of Christ is a sin of extreme gravity,'' and
``with a sorrowful yet hopeful heart, I feel it is my duty, through the
authority received from Christ, to ask you to desist from your intended
act.''\endnote{Leo XIV, letter to the Reverend Father Davide Pagliarani,
Superior General of the Priestly Fraternity of Saint Pius X, 29 June
2026; French original and official English on the Holy See website
(\url{https://www.vatican.va/content/leo-xiv/en/letters/2026/documents/20260629-lettera-fraternita-sanpiox.html}),
published in the Holy See Press Office bulletin B0563 of 30 June 2026.
The block quotations are from the Holy See's own English version. Read
2026-07-25.}

\actrecord{Letter of Leo XIV, 29 June 2026 (French original; official
English on the Holy See website).}{A personal papal request, made the day
before the ceremony and by the authority received from Christ, that the
Society desist; a papal characterization of the intended act as
schismatic; and a papal warning that it would deprive the faithful of the
licit and in some cases valid reception of the sacraments.}{An appeal, not
a penal decree. It imposes nothing and declares nothing; the penal act
came on 2 July from the doctrinal dicastery. Its warning about validity
is stated in general terms and is not, in the letter itself, tied to
named sacraments or to the revocation of any earlier faculty.}

\subsection{1 July 2026: the consecrations}

On 1 July 2026, the feast of the Most Precious Blood, at the seminary of
Saint Pius X at Écône, Bishop Alfonso de Galarreta, assisted by Bishop
Bernard Fellay, consecrated the four priests. The Society's statement of
that day describes them as consecrated ``to serve as auxiliary bishops of
the Society of Saint Pius X, without jurisdiction,'' regrets that the
consecrations ``had to be conferred without the authorization of the Holy
Father'' and that the Superior General had not been received personally
by Leo XIV, and gives thanks for what it calls a very great grace for the
Society and for the whole Church.\endnote{``General House Statement
Following the Episcopal Consecrations,'' Écône, 1 July 2026
(fsspx.news), read 2026-07-25. Party statement. The Society reports a
very great gathering of priests, religious and faithful; its later
timeline gives a figure of 17{,}000 faithful from nearly seventy
countries, which is self-reported and not verified here.}

At the point in the rite where the apostolic mandate is read, the
Society's Secretary General instead read a declaration setting out the
reasons for proceeding. Its published English text states that ``entirely
exceptional circumstances'' require the Society to provide for the
transmission of the deposit of faith; that since the Second Vatican
Council ``the authorities in the Church have been animated by a spirit
that is contrary to that of the Faith''; that it is therefore a sacred
duty before God to consecrate bishops faithful to Tradition; and---the
sentence that defines the canonical dispute---``we consider that every
punishment and censure brought to bear against this step will have no
validity.''\endnote{``\guillemotleft Habetis mandatum apostolicum?
\guillemotright{} \textbar{} Declaration read before the episcopal
consecrations'' (fsspx.news, 2 July 2026), read 2026-07-25. Party text.
The declaration also names the four priests with their offices.}

\subsection{2 July 2026: the decree}

On 2 July 2026 the Dicastery for the Doctrine of the Faith published a
decree and an explanatory note, both under protocol number 99/2009---the
same file number the Roman dossier has carried since the year of the
remission. The decree is signed by Cardinal Víctor M. Fernández, Prefect;
Archbishop John J. Kennedy, titular of Ossero, Secretary for the
Disciplinary Section; and Monsignor Armando Matteo, Secretary for the
Doctrinal Section. Its content, in order:\endnote{Dicastery for the
Doctrine of the Faith, \work{Decreto}, Prot.\ N.~99/2009, 2 July 2026,
Holy See Press Office bulletin B0568
(\url{https://press.vatican.va/content/salastampa/it/bollettino/pubblico/2026/07/02/0568/01077.html}),
read 2026-07-25. Italian original; the summary below paraphrases it
clause by clause and offers no project translation.}

\begin{enumerate}
\item Notwithstanding the admonitions addressed to the Superior General
of the Society, Bishop Alfonso de Galarreta, having performed an act of a
schismatic nature by the episcopal consecration of four presbyters
without pontifical mandate and against the will of the Supreme Pontiff,
incurred \latin{ipso facto} the penalties provided by canon 1387 and
canon 1364 \S\,1 of the Code as revised in 2021.
\item The decree declares, for all juridical effects, that the said
Bishop de Galarreta and Pascal Schreiber, Michael Goldade, Michel
Poinsinet de Sivry and Marc Hanappier have incurred \latin{ipso facto}
the excommunication \latin{latae sententiae} reserved to the Apostolic
See.
\item It declares further that Bishop Bernard Fellay, having taken part
directly in the liturgical celebration as co-consecrator and having thus
publicly adhered to the schismatic act, incurred the excommunication
\latin{latae sententiae} provided by canon 1364 \S\,1.
\item It admonishes clerics and lay faithful not to adhere to the schism
of the Society, since they would \latin{ipso facto} incur the penalty of
excommunication \latin{latae sententiae}.
\end{enumerate}

Two differences from 1988 are worth marking. The consecrator is
penalized under both the consecration canon and the schism canon, and the
co-consecrator under the schism canon by reason of public adherence---a
route not used in the 1988 declaration, which rested on the consecration
canon alone. And the final admonition is addressed to clerics and laity
generally, in advance.

\subsection{The explanatory note}

The note, published with the decree and under the same protocol number,
is the document with the wider reach. It opens by observing that from the
time of Saint Paul VI to the most recent conversations held at the
Dicastery, the many attempts to bring the adherents of the movement begun
by Archbishop Lefebvre to full communion have proved vain, and that the
situation has been further aggravated by the recent consecrations; the
Dicastery therefore judges it necessary to note that the act has
constituted the delict of schism, with canonical consequences for sacred
ministers and lay faithful, recalling \work{Ecclesia Dei} n.~3 on
disobedience that carries a practical rejection of the Roman primacy.
Then, ``from now on'':\endnote{Dicastery for the Doctrine of the Faith,
\work{Nota Esplicativa}, Prot.\ N.~99/2009, 2 July 2026, same bulletin
B0568
(\url{https://press.vatican.va/content/salastampa/it/bollettino/pubblico/2026/07/02/0568/01078.html}),
read 2026-07-25. Italian original; the summary paraphrases it. The note
cites \work{Ecclesia Dei} 3 and 5~c), the 1996 note of the Pontifical
Council for Legislative Texts at nn.~5--7, canon 1364 \S\,1, canon 751,
and \work{Lumen gentium} 22.}

\begin{itemize}
\item the sacred ministers belonging to the Society ``are in schism and
must therefore be considered schismatics,'' and are subject to the
excommunication provided by law in canon 1364 \S\,1;
\item lay faithful who formally adhere to the Society under the
conditions established in the 1996 note of the Pontifical Council for
Legislative Texts---``still in force, which this Dicastery makes its
own''---are to be held schismatic and excommunicated;
\item the holy People of God is warned that the Society's sacred
ministers administer the sacraments illicitly, and that ``the sacrament
of penance administered by them and the marriage at which they assist are
invalid'';
\item the Church will welcome with sincere affection those who wish to
return to full communion, and the apostolic nuncios will make available
the procedures that ordinaries may use in the various cases;
\item all the faithful are exhorted to remain firm in communion with the
Roman Pontiff, with the bishops in communion with him and with the whole
Church, and to abstain from taking part in the celebrations and
activities promoted by the Society.
\end{itemize}

The note thereby does two things the decree does not. It adopts the 1996
consultative note as the Dicastery's own---which converts a curial
council's opinion into the operative standard of a competent dicastery,
including its case-by-case test for lay adherence. And it states the
sacramental consequences that follow when a minister has no faculty:
penance and marriage, the two sacraments whose validity depends on
faculty or delegation, are declared invalid, while the note conspicuously
does not say the same of the Eucharist, confirmation or orders, which it
describes only as illicitly administered.

\subsection{The Society's response and the recourse}

The Society answered on three levels. On 3 July the Superior General
wrote to Leo XIV, in a letter built on the Gospel figures of the stone,
the serpent and the scorpion, saying that the Society had asked for
understanding and had received condemnation, promising nonetheless that
it would not receive ``these new sanctions---objectively unjust and
invalid---with bitterness or revolt.''\endnote{``Letter to the Holy
Father regarding the Decree of the Dicastery for the Doctrine of the
Faith,'' Écône, 3 July 2026 (fsspx.news), read 2026-07-25. Party
document; quotations are from the Society's own English text.}

On 11 July it lodged a canonical recourse, announced in a communiqué of
13 July: in response to the decree of 2 July, the Society ``submitted on
11 July a preliminary recourse to the same Dicastery, in accordance with
canons 1734 and following of the Code of Canon Law,'' this being ``the
mandatory preliminary step before the possible introduction of a
hierarchical recourse,'' and having, in the Society's words, ``the effect
of suspending the execution of the decree, in accordance with canon
1353.''\endnote{``General House Communiqué: The Society Files a Recourse
Against the Decree of 2 July 2026,'' Menzingen, 13 July 2026
(fsspx.news), read 2026-07-25. Party document.}

The canons named are real and say what the Society says they say, up to a
point that the Society's communiqué passes over. Canon 1734 \S\,1 does
require a petition for revocation or emendation to the author of a decree
before recourse, and provides that the making of that petition is
understood to include a request for suspension of execution; canon 1353
does provide that appeal or recourse from judicial sentences or from
decrees which impose or declare any penalty has suspensive effect.
Whether a declaratory decree concerning penalties incurred \latin{ipso
facto} produces, through that canon, the suspension the Society claims,
and what such a suspension would mean for penalties already incurred by
the law itself, is a contested question of canon law. The Holy See has
published nothing on it. This account records the claim, records its
legal basis, and does not decide it.\endnote{Canon 1353 (\work{Codex
Iuris Canonici}, Book VI as revised in 2021): ``Appellatio vel recursus a
sententiis iudicialibus vel a decretis, quae poenam quamlibet irrogent
vel declarent, habent effectum suspensivum''; canon 1734 \S\S\,1--2 and
canon 1737 (Book VII) on the preliminary petition, the ten useful days
for it, and the fifteen useful days for hierarchical recourse. Latin
texts read in the Holy See's official web presentation of the Code on
2026-07-25 and registered in this repository's source library. The
canons are quoted for what they provide; their application to this
decree is precisely what is disputed.}

Finally, on 19 July, the Superior General wrote to the faithful,
describing the ceremony and the storm that broke over it, calling the
condemnation ``wholly unjust'' and an ``immense humiliation,'' telling
the faithful that ``it is Tradition itself \ldots that is excommunicated
in all those who embody it,'' and that ``before God, such a condemnation
cannot be valid.'' The letter contains no announcement of any change of
practice.\endnote{``Letter from the Superior General to all the faithful
of the Society of Saint Pius X,'' Menzingen, 19 July 2026 (fsspx.news),
read 2026-07-25. Party document; quotations from the Society's own
English text.}

\subsection{Where the file stood at the end of July 2026}

As of 25 July 2026 the Holy See had published no response to the
recourse. The Holy See Press Office bulletins from 2 to 25 July 2026
contain no further item concerning the Society, and no act of the
Dicastery, the Secretariat of State, or the pope on the matter appeared
on the Holy See website after 2 July.\endnote{Bounded negative: every
bulletin listing for July 2026 was checked on 2026-07-25, and the only
Society-related items are the decree and note of 2 July (bulletin B0568).
This is a bounded search of one publication channel on one date; a
non-public act, or an act published elsewhere, would not appear in it.}
